\chapter{Элементарный объём и материя}
\label{ch:matter}

\section{Бинарная занятость и первичный дефект}

На каждой ячейке~$\PlanckCell$ вводится бинарная занятость $b(x)\in\{0,1\}$:
фон вакуума ($b=0$) или локализованный узел ($b=1$).
Плотность материи принимает значения $0$ или планковская~$\rho_P$.
Первичный дефект --- одна ячейка, а не «облако'' внутри~$\hl$.

\section{Несжимаемая решёточная жидкость}

\begin{definition}[Модуль упругости $K_P$]
\label{def:kp}
$K_P=u_P$ --- планковское давление жидкости;
вращатель~$\Phi$ насыщается при $\rho\to u_P$.
\end{definition}

Отклонение плотности от фона вызывает упругий фазовый отклик;
избыток энергии уходит волной ($n=0$) или вихрем.
Уравнения Эйлера и Навье--Стокса на~$T$ --- макроосреднение,
а не буквальное дискретное уравнение на~$M$.

\section{Теорема о механической лестнице}
\label{sec:theorem-51}

\begin{theorem}[Лестница Ландау на ячейке]
\label{thm:51}
После \cref{thm:discreteness,thm:g-law} локальные механические величины на~$\hv$ образуют
единственную лестницу $s_0\to p_0\to L_0\to F_0\to E_0$
без трёх независимых подгоняемых параметров.
\end{theorem}

\begin{lemma}[5.1a]
\label{lem:51a}
Из~$A_5$ и $\Delta\phi_{\min}=\tfrac12$ следует $s_0=\hbar/2$.
\end{lemma}

\begin{proof}
Минимальное действие одного события изменения аргумента на~$\hv$ за такт~$\htick$.
\end{proof}

\begin{lemma}[5.1b]
\label{lem:51b}
Из $s_0$ и~$\hl$: $p_0=s_0/\hl=\hbar/(2\hl)$, $E_0=s_0/\htick$, $m_{\arg}=E_0/c_0^2$.
\end{lemma}

\begin{lemma}[5.1c]
\label{lem:51c}
$L_0=p_0\hl=s_0$, $g_M=c_0/(2\htick)$, $F_0=p_0/\htick=m_{\arg} g_M$ ---
соотношение $F=mg$ на одной ячейке~$\hv$.
\end{lemma}

\begin{proof}
Арифметика $(\hl,\htick)$ без дополнительных констант.
\end{proof}

\begin{lemma}[5.1d]
\label{lem:51d}
Ударный член $\lfloor\mathcal{N}\rfloor$ на~$\mathbb{Z}$ $\Rightarrow$ $\Delta p\in p_0\cdot\mathbb{Z}$;
интеграл $\int\Im(z^*\nabla z)$ --- только макроописание на~$T$.
\end{lemma}

\begin{lemma}[5.1e]
\label{lem:51e}
$E_0=p_0 c_0=F_0\hl=L_0/\htick$, $E=n_E E_0$, $n_E\in\mathbb{Z}$.
\end{lemma}

\begin{proof}[Доказательство теоремы]
Следует из \cref{lem:51a,lem:51b,lem:51c,lem:51d,lem:51e}.
\end{proof}

\begin{corollary}[Закон Ньютона на ячейке]
$p_0=m_{\arg}(c_0/2)$, $F_0=m_{\arg} g_M$, $\Delta F=\Delta m\cdot g$ при $\Delta m\in m_{\arg}\mathbb{Z}$.
\end{corollary}

\section{Симметрии и дискретные заряды}

\begin{proposition}[Теорема Нётер на решётке]
На~$M$ сохраняются:
\begin{itemize}
\item заряд при $U(1)$: $Q=\sum|z|^2$;
\item импульс: $\mathrm{div}\,\pi=0$ по модулю~$p_0$;
\item момент при $SO(2)$ изотropного~$N$: $L_z\in L_0\cdot\mathbb{Z}$;
\item энергия: $E\in E_0\cdot\mathbb{Z}$;
\item топологический заряд: $n\in\mathbb{Z}$ ($A_{10}$).
\end{itemize}
\end{proposition}

\section{Геометрический квант связи \texorpdfstring{$\kappa_{\mathrm{link}}=1/|N|$}{kappa=1/|N|}}

\begin{proposition}
На FCC $|N|=12$ $\Rightarrow$ $\kappa_{\mathrm{link}}=1/12=\gamma=\nu_{\mathrm{CA}}/(c_0\hl)$.
На гексагональном срезе $|N|=6$ $\Rightarrow$ $1/6$.
\end{proposition}

Спектральные методы калибруют макроуровень~$T$, а не задают закон на~$M$.

\section{Фононы и спектр}

Периодичность~$\Lambda$ $\Rightarrow$ первая зона Бrillouin;
$p=n p_0+\hbar k$, $k$ дискретен.
Фонон вакуума: $\omega(k)\approx v_a|k|$, $v_a=2c_0$.
Фонон вещества --- второй кристалл над микроуровнем (гл.~\ref{ch:carrier}).
